\clearpage\fancyhead[L]{\small JUNO.}\hypertarget{native-365-lower}{}\label{native-365-lower}
\begin{center}\textbf{JUNO.}\par\vspace{4pt}\hrule\vspace{4pt}\fontsize{10}{12}\selectfont Monatliche Correspondenz. Band XXVI. Seite 297\;..\;299. 1812 September.\par\vspace{4pt}\hrule\vspace{4pt}\end{center}

\begin{flushright}Göttingen den 9. Sept. 1812.\end{flushright}
\par\indent Heute will ich meine \emph{Juno}-Beobachtungen nachholen. Das Auffinden
dieses kleinen Planeten machte dieses Jahr überaus viel Mühe, da er
kaum die 11te Grösse hatte und gerade in der Milchstrasse stand. Dazu kam
noch, dass die Aufsuchung durch Mondschein und schlechtes Wetter
unterbrochen wurde. So habe ich nur vier Beobachtungen erhalten, die ich
aber alle für sehr gut halte.

\begin{center}\fontsize{10}{12}\selectfont\setlength{\tabcolsep}{4pt}\renewcommand{\arraystretch}{1.15}\begin{tabular}{l|c|c|c}
1812 & Mittlere Zeit & AR. & Declination\\\hline
Juni 7 & $12^{\mathrm u}\;21^{\mathrm m}\;0^{\mathrm s}$ & $282^\circ\;21'\;24''5$ & $4^\circ\;41'\;44''6$ südl.\\
13 & $12\;51\;12$ & $281\;16\;32.9$ & $4\;33\;39.6$\\
Juli 12 & $11\;35\;28$ & $275\;5\;10.3$ & $5\;6\;11.2$\\
13 & $11\;28\;12$ & $274\;52\;49.9$ & $5\;9\;20.0$\\
\end{tabular}\end{center}

\Needspace{13\baselineskip}\par\indent Meine X Elemente geben folgende Unterschiede:

\begin{center}\fontsize{10}{12}\selectfont\setlength{\tabcolsep}{7pt}\renewcommand{\arraystretch}{1.15}\begin{tabular}{l|c|c}
& AR. & Declinat.\\\hline
Juni 7 & $+21''8$ & $-35''9$\\
13 & $+31.0$ & $-34.7$\\
Juli 12 & $+22\;4$ & $-46.9$\\
13 & $+26.2$ & $-45.0$\\
\end{tabular}\end{center}

\par\indent Herr \textsc{Wachter} übernahm die Arbeit, aus diesen Beobachtungen die Opposition
zu rechnen. Ich empfahl ihm hierzu eine Methode, welcher ich mich schon
seit langer Zeit zu dieser Absicht bedient habe, und die den Vortheil
gewährt, dass man dabei ohne Bedenken auch etwas entfernter liegende
Beobachtungen für die Opposition benutzen kann; ich werde Ihnen ein
andermal etwas umständlicher über dieses Verfahren schreiben. Herrn
\textsc{Wachter}'s Resultat ist folgendes:

\Needspace{8\baselineskip}\begin{center}
\emph{Opposition der Juno} 1812.
\end{center}

den 29. Juni $19^{\mathrm u}\;51^{\mathrm m}\;43^{\mathrm s}$ Mittlere Zeit in Göttingen

\begin{center}\fontsize{10}{12}\selectfont\setlength{\tabcolsep}{2pt}\renewcommand{\arraystretch}{1.15}\begin{tabular}{lcr}
wahre Länge & $.\;.\;.\;.\;.\;.\;.\;.\;.\;.\;.$ & $278^\circ\;15'\;21''3$\\
wahre geocentrische Breite & $.\;.\;.\;.\;.\;.$ & $18\;36\;56.9$\\
\end{tabular}\end{center}

\Needspace{19\baselineskip}\par\indent Herr \textsc{Wachter} hat ferner aus den vier letzten beobachteten Oppositionen,
(die von 1810 ist bekanntlich versäumt) die Elemente corrigirt, nach der
im I. Bd. der \emph{Comment. Nov. Soc. Sc. Götting.} [S. 1. d. B.]
vorgetragenen Methode, und gefunden:
\SourceLeaf{366}\begin{center}\fontsize{9}{11}\selectfont\setlength{\tabcolsep}{2pt}\renewcommand{\arraystretch}{1.15}\begin{tabular}{lcr}
Epoche 1813 Meridian von Göttingen & $.\;.\;.\;.\;.\;.\;.$ & $342^\circ\;56'\;32''2$\\
Tägliche mittlere tropische Bewegung & $.\;.\;.\;.\;.\;.$ & $813''25748$\\
Perihelium 1813 & $.\;.\;.\;.\;.\;.\;.\;.\;.\;.\;.\;.\;.\;.$ & $53^\circ\;16'\;52''3$\\
Knoten 1813 & $.\;.\;.\;.\;.\;.\;.\;.\;.\;.\;.\;.\;.\;.\;.\;.$ & $171\;10\;58.9$\\
Neigung der Bahn & $.\;.\;.\;.\;.\;.\;.\;.\;.\;.\;.\;.\;.$ & $13\;4\;17.2$\\
Excentricitäts-Winkel & $.\;.\;.\;.\;.\;.\;.\;.\;.\;.\;.\;.$ & $14\;43\;58.1$\\
Logarithm. der halben grossen Axe & $.\;.\;.\;.\;.\;.\;.$ & $0.4265679$\\
\end{tabular}\end{center}

\par\indent Da jetzt die Säcularänderungen der Elemente der Ceres- und Pallasbahn, so weit sie vom Jupiter
herrühren (gegen dessen Wirkung die übrigen Planeten ganz unbedeutend sind) bekannt sind, so
veranlasste ich Herrn \textsc{Enke}, einen talentvollen jungen Mann und eben so geschickten als sorgfältigen
Rechner, zu untersuchen, ob der Abstand der Ceres- und Pallasbahn, der bekanntlich beim aufsteigenden
Knoten der Ceresbahn auf der Pallasbahn ziemlich klein ist, — wodurch eben Herr Dr. \textsc{Olbers} auf seine
bekannte Hypothese geleitet wurde — jetzt im Abnehmen oder Zunehmen sei, d.\,i. ob die Bahnen auf einen
wirklichen Schnitt zugehen, oder davon herkommen. Das letztere würde offenbar \textsc{Olbers} Hypothese günstig
sein, allein das Resultat ist gerade umgekehrt. Die Distanz ist ehemals grösser gewesen, als sie jetzt
ist. Herr \textsc{Enke} findet die Radios vectores

\begin{center}\fontsize{9}{11}\selectfont\setlength{\tabcolsep}{2pt}\renewcommand{\arraystretch}{1.15}\begin{tabular}{c|c|c|c||c|c|c}
& \multicolumn{3}{c||}{\shortstack{im $\AscendingNode$ der Ceresbahn auf der\\Pallasbahn}} & \multicolumn{3}{c}{\shortstack{im $\DescendingNode$ der Ceresbahn auf der\\Pallasbahn}}\\
Jahr & Ceres & Pallas & Unterschied & Ceres & Pallas & Unterschied\\\hline
$808$ & $2.82294$ & $2.70322$ & $-0.11972$ & $2.67612$ & $2.37780$ & $-0.29832$\\
$1808$ & $2.92427$ & $2.84945$ & $-0.07482$ & $2.59569$ & $2.40346$ & $-0.19223$\\
$3475$ & $2.95374$ & $2.95743$ & $+0.00369$ & $2.57987$ & $2.49163$ & $-0.08824$\\
\end{tabular}\end{center}

\par\indent Hiernach würde etwa um das Jahr 3397 ein wirklicher Schnitt im $\AscendingNode$ erfolgen, welches man immer
als näherungsweise richtig betrachten darf. Freilich wird auch einmal ein Schnitt Statt gefunden haben;
allein aus dem Gange der Zahlen in der vierten und sechsten Columne lässt sich wenigstens schliessen,
dass dies nur zu einer viele Jahrtausende entfernten Epoche möglich gewesen sein könne. Wenn man also
Dr. \textsc{Olbers} Hypothese über den Ursprung der neuen Planeten annehmen will, so fällt derselbe in eine
noch gar nicht zu berechnende Ferne vor die Zeiten, wohin unsre Geschichte reicht.
